\section{Convergence of correlations along the full mesh sequence}
\label{sec:fullmesh-correlation-cauchy}
We now pass from convergence of the local data at each fixed level to
convergence of the finite-volume correlations themselves. The two limits
are different: a terminal correlation contains interactions involving
several marked insertions, whereas the fixed-level convergence theorem
controls the single-insertion coefficients and their forced profiles.
The missing estimate makes the effect of the former interactions small
at a sufficiently fine intermediate level, uniformly in the terminal
cutoff.

Fix a period $p$, a finite set $V$ of marked insertions, their grades
$d_i$, coefficient vectors $z_i$, and one of the modes $S$ of their
finite marker algebra. Throughout the argument the history, the local
fields, and the transfer maps belong to the one selected family.
At cutoff $N$, let $R_{N,m}$ be its representation at level $m\le N$.
Its product-Haar prior is denoted by $\mu_{m,p}$, and the positive
terminal Wilson partition function by $Z_{N,p}$. The transferred
polynomial density $X_{N,m}$ satisfies the exact identities
\[
 X_{N,m}=\mathcal T_{N,m}X_{N,m+1},\qquad
 \int X_{N,m}\,d\mu_{m,p}
       =\int X_{N,m+1}\,d\mu_{m+1,p}.
\]
Here $\mathcal T_{N,m}$ is the actual gauge-covariant transfer at the
coupling of $R_{N,m+1}$; it is applied coefficient by coefficient.
The integral of the coefficient of mode $S$, divided by $Z_{N,p}$, is the correlation
under consideration. In particular all normalizations in this section
use the same partition function.

\subsection{The fixed-level data and their uniform bounds}
At level $m$ write $\Phi_{N,m,i}(U)$ for the single-insertion term in
the exponent and $\omega_{N,m,i}(Q,U)$ for its forced profile, where
$Q$ records a support together with its exact set of flagged sites.
Both include the field's scaling factor, its translated origin, and
its prescribed motion. These are the actual transferred fields of
Section~\ref{sec:selected-history-limits}.

For every fixed $m$ above the entry levels of the finitely many grades,
the fixed-level comparison supplies uniform Cauchy convergence in $U$
of each $\Phi_{N,m,i}$ and each $\omega_{N,m,i}(Q,\cdot)$, and
eventual uniform boundedness. This assertion uses both parts of that
comparison. The finite-dimensional principal coefficient vectors are
Cauchy. The captured remainder vectors are Cauchy in the local norm
which controls their evaluation uniformly in $U$. Evaluation, translation,
and the fixed motion preserve these bounds. The local masks are the
same at this fixed level for every sufficiently large $N$, so
multiplying by a mask preserves uniform Cauchy convergence even if
the mask is discontinuous. The selected paths agree with the stationary
paths of the common horizon; their larger-width convergence supplies
the path part of the vacuum comparison. The fixed-level activity
comparison supplies its remaining part.

Here and below finite initial segments cause no difficulty. A Cauchy
sequence of coefficient vectors is bounded, including its finite initial
segment. Since the coefficient price is a continuous seminorm times
the fixed scaling factor, there is a number $A_i$ such that the actual
paid coefficient of insertion $i$ at a fixed root $r$ is at most $A_i$
for every cutoff. Choose
\begin{equation}\label{eq:fullmesh-root-weights}
 w_i=\frac{\zeta_{d_i}(r)}{|A_i|+1}>0,
\end{equation}
where $\zeta_d(r)>0$ is the available single-insertion allowance.
Then $w_i$ times that paid coefficient is at most $\zeta_{d_i}(r)$.
The root $r$ is chosen after the finite list of grades, so these
allowances are positive there.

For later use, at any fixed level $m$ and for any fixed positive
weights $v_i$, there exist $x_m,y_m\ge0$, independent of all
sufficiently large $N$, such that
\begin{equation}\label{eq:fullmesh-vault}
 v_i|\Phi_{N,m,i}(U)|\le x_m,\qquad
 \sum_Q v_i|\omega_{N,m,i}(Q,U)|\le y_m
 \quad(i\in V).
\end{equation}
Indeed, take eventual bounds $a_i$ and $b_{i,Q}$ for these functions
and put $x_m=\sum_i v_i|a_i|$ and
$y_m=\sum_{i,Q}v_i|b_{i,Q}|$. There are only finitely many supports
and flag inventories on the fixed torus. Taking the maximum of the
finitely many entry thresholds proves both inequalities simultaneously.
Measurability follows from the continuous local functions, measurable
masks and absolutely summable forced series in their construction.

\subsection{Separating the higher marker terms}
Write $\mathsf M(\Phi)=\sum_i\Phi_i e_{\{i\}}$ for the
single-marker polynomial. With the level and cutoff suppressed, the
transferred density has the exact form
\begin{equation}\label{eq:fullmesh-restore}
 \begin{split}
 \mathcal R(H,K)(U)={}&e^{R.\mathrm{exponent}(U)}
 \operatorname{Pass}\!\left(\mathsf M(\Phi(U))+\sum_P H_P(U)\right)
 \\ &{}\times \mathcal F\!\left(\overline R_Q(U)+\mathsf M(\omega(Q,U))+K_Q(U);
                         \mathrm{flags}(U)\right).
 \end{split}
\end{equation}
The finite nilpotent exponential $\operatorname{Pass}$ was defined in
Section~\ref{ex-centered-construction}. In this formula
$\overline R_Q$ is the vacuum forced profile. For an inventory $F$,
$\mathcal F(A;F)$ is the sum of $\prod_{Q\in\mathcal Q}A_Q$ over
families $\mathcal Q$ with disjoint supports whose flag inventories
have union $F$. This is a finite sum on the fixed torus.
Every nonzero mode in $H_P$ or $K_Q$ contains at least two markers.
Its marker origins lie in its support. The functions $H_P$ vanish
off the healthy set of $P$; the functions $K_Q$ vanish unless the
flag inventory is exactly that specified by $Q$. Both are local,
gauge-invariant and measurable.

For positive marker weights $v_i$, use the weighted coefficient price
of Section~\ref{ex-centered-construction}. The higher terms have
nonnegative majorants $a_P$ and $b_Q$ bounding this price uniformly
in the field. Their rooted support budgets, with weight
$e^{e_m\operatorname{cov}_{B_m}(P)}$, are at most one, for both
$H$ and $K$. Here $e_m$ is the level's support exponent,
$B_m$ its support scale and $\operatorname{cov}_{B_m}$ the cover
number already used in those budgets.

We explain why these statements concern the actual density. The
centered terminal construction gives \eqref{eq:fullmesh-restore} at
level $N$ with exponent $e_N-3$. Its price uses
\[
 c_N=\kappa_N+r_N\operatorname{siteFactor}(B_N)u_N,
 \qquad (2B_N+1)^4c_N\le3,\qquad e^{c_N}-1\le1,
\]
where $r_N$ is the allowed number of origins per site, $\kappa_N$
the vacuum forced allowance and $u_N$ the terminal insertion cost.
These are the terminal estimates in the centered birth construction.
At a step $m+1\to m$, the purchased finite-menu estimate accepts
support exponent $e_{m+1}-4$ and produces the full exponent $e_m$.
The incoming exponent is sufficient, including at $m+1=N$, because
$e_N-4\le e_N-3$. The one-insertion transfer identities identify
its single-marker coefficients with the actual $\Phi_{N,m,i}$ and
$\omega_{N,m,i}$; they are not replaced by estimates. The local
transfer ledger and the occupancy budget then give precisely the
higher-term majorants above. Integrability propagates through the
same transfer. Backward induction proves the assertion at every
$m<N$.

The hypotheses for this induction are the funded finite-menu service
of Section~\ref{sec:finite-menu-history-handoff} and the centered
service of Section~\ref{ex-centered-public-service}. Their
single-insertion component is obtained as follows. At the root use
\eqref{eq:fullmesh-root-weights}, and propagate weights by
$v_i(m+1)=\varphi_m v_i(m)$. The centered scale estimates give
\[
 0<\varphi_m,\quad 0<\delta_m,\quad
 1\le\delta_m\varphi_m e^{B_m/128},\quad
 B_{m+1}^2\le L_m B_m^2,\quad B_m\ge8,\quad e_m\ge4.
\]
The actual single-insertion transfer estimate and its coefficient
price recursion propagate the root inequality to every later level.
The terminal local section has cover number at most one and its
purchased coefficient bound. The coefficient generators in that
estimate therefore give the terminal insertion, the exact
single-insertion recursion, and its local finite transfer ledgers.
Summing their forced coefficients by support and exact flag inventory
returns the original $\omega_{N,m,i}$: this is the defining profile
identity, not a change of the fields. The remaining applicability
conditions are the funded common service and the convergence
comparisons just described; their simultaneous production is required
in the common-family assembly.

\subsection{Removing the higher terms at a fine intermediate level}
Suppose distinct origins at level $j$ are separated by at least $B_j^2$
in the lattice metric. A support of a nonzero higher term contains
at least two such origins. The cover estimate of
Section~\ref{ex-centered-construction} gives
\[
 B_j\le64\operatorname{cov}_{B_j}(P).
\]
Put $S_j=e^{B_j/64}$. For $|z|\le S_j$, replace $H,K$ by $zH,zK$.
The support and flag conditions are unchanged. The price majorant
of a nonzero term is multiplied by $|z|$, and
\begin{equation}\label{eq:fullmesh-inflation}
 |z|e^{(e_j-2)\operatorname{cov}_{B_j}(P)}
 \le e^{\operatorname{cov}_{B_j}(P)}
           e^{(e_j-2)\operatorname{cov}_{B_j}(P)}
 \le e^{e_j\operatorname{cov}_{B_j}(P)}.
\end{equation}
Zero terms retain the zero majorant. Thus the inflated density has
the same cost one at exponent $e_j-2$. Its single-insertion terms
are unchanged. Transport it from $j$ down to a fixed root $r<j$.
The first transfer accepts it since $e_j-4\le e_j-2$; subsequent
steps again have the full exponent. The preceding induction therefore
applies, with the same root single-insertion terms and weights.

For clarity, separation also supplies the weight reservation needed
in this transport. If a mode contains $q\ge2$ distinct origins, the
stronger cover estimate is $B_m q\le64\operatorname{cov}_{B_m}(P)$.
Changing each marker weight by at most $e^{B_m/128}$ multiplies its
toll by at most $e^{\operatorname{cov}_{B_m}(P)/2}$. This pays the
half-unit reservation in the finite-menu transfer. The identities
between parent origins preserve occupancy, and
$B_{m+1}^2\le L_mB_m^2$ preserves separation under the exact
lattice embeddings. These observations verify the geometric
conditions at every step of the transport.

Let $X_{N,j}=\mathcal R(H,K)$ be the actual density at $j<N$, and
set
\[
 F_N(z)=\frac1{Z_{N,p}}
   \int\bigl(\mathcal R(zH,zK)(U)\bigr)_S\,d\mu_{j,p}(U).
\]
This is holomorphic in $z$. Indeed, the marker algebra is nilpotent
and the inventory sum finite, so the expression is a polynomial in
$z$ with measurable coefficients. The price bounds give integrable
coefficient majorants on every smaller disk. At the root, let
$x_r,y_r$ be the eventual bounds \eqref{eq:fullmesh-vault} with
weights $w_i$. The normalized density estimate proved in
Section~\ref{ex-centered-construction} gives
\begin{equation}\label{eq:fullmesh-root-bound}
 |F_N(z)|\le C_r:=
 \frac{\exp\{ |\Lambda_{r,p}|\tau_r+|V|(x_r+y_r+2)\}}
 {\operatorname{toll}_w(S)}\qquad(|z|<S_j).
\end{equation}
To apply it, transfer the inflated density to $r$. Its integral and
its vacuum normalization remain exactly those above. The support
budgets at $r$ are one, and its single-insertion data have the root
bounds $x_r,y_r$. This proves \eqref{eq:fullmesh-root-bound} with
a constant independent of $j$ and of all sufficiently large $N$.
The bounds at $j$ needed for integrability may depend on $j$; they
do not enter $C_r$.

If $S_j>2$, Cauchy's estimate bounds $|F_N'(z)|$ by $2C_r/S_j$
for $|z|\le1$: the closed disk of radius $S_j/2$ about such a
point lies in $|z|<S_j$. Integrating the derivative on $[0,1]$
therefore gives
\begin{equation}\label{eq:fullmesh-omission}
 \left|\frac1{Z_{N,p}}\int(X_{N,j})_S\,d\mu_{j,p}
       -\frac1{Z_{N,p}}\int(\mathcal R(0,0))_S\,d\mu_{j,p}\right|
 \le2C_r e^{-B_j/64}.
\end{equation}
The first term is the original terminal correlation. Only the terms
with at least two markers were erased in the second term.

\subsection{Cauchy convergence of the erased expression}
Fix $j$. Remove the scalar normalization from the representation and
put
\[
 D_{N,j}(U)=e^{-R_{N,j}.\mathrm{scalar}}\mathcal R(0,0)(U),
 \qquad
 \ell_{N,j}=e^{-R_{N,j}.\mathrm{scalar}}Z_{N,p}.
\]
The second term in \eqref{eq:fullmesh-omission} is exactly
$\int(D_{N,j})_S\,d\mu_{j,p}/\ell_{N,j}$; the exponential
scalar cancels. The fixed-level vacuum estimate gives a number
$a_{j,p}>0$ with $\ell_{N,j}\ge a_{j,p}$ eventually, and the
fixed-level convergence theorem gives Cauchy convergence of
$\ell_{N,j}$.

The functions $(D_{N,j})_S$ are uniformly Cauchy and eventually
uniformly bounded. To see this directly from
\eqref{eq:fullmesh-restore}, combine the fixed-level vacuum, activity,
single-insertion and forced-profile comparisons. There are finitely
many coordinates on the fixed torus. Finite products, sums and the
nilpotent exponential are continuous functions of those coordinates,
uniformly continuous on the compact closures of their bounded ranges.
Fixed masks can be kept in every term, as noted earlier. These
operations consequently preserve uniform Cauchy convergence.
As the prior is a probability measure,
\[
 \left|\int(D_{N,j})_S\,d\mu_{j,p}
            -\int(D_{M,j})_S\,d\mu_{j,p}\right|
 \le\sup_U|(D_{N,j}(U)-D_{M,j}(U))_S|\longrightarrow0.
\]
The numerator is Cauchy and the denominator has a positive limit,
so their ratio is Cauchy as well.

Now $B_j\to\infty$. Given $\varepsilon>0$, first choose a fixed
$j>r$ with $S_j>2$ and $2C_re^{-B_j/64}<\varepsilon/4$.
At this fixed $j$ the erased expressions differ by less than
$\varepsilon/4$ for sufficiently large cutoffs $N,M>j$.
Equation~\eqref{eq:fullmesh-omission} and the triangle inequality
bound the difference of the original correlations by
$3\varepsilon/4<\varepsilon$. This proves Cauchy convergence along
the full cutoff sequence, without choosing an additional subsequence.

\subsection{From nested-grid points to all distinct configurations}
It remains to supply the separated origins used above and to pass
to general points. Call a physical point ancestral if it lies on the
mesh at some level. Integer block lengths imply that it lies on every
finer mesh, and its addresses obey the exact embedding identity.
Ancestral configurations are dense: round each coordinate to the
nearest mesh point and use $a_m\to0$.

For a fixed ancestral configuration of distinct points on the period
$p$ torus, let $\eta>0$ be its minimum pairwise physical separation.
For zero or one insertion take any $\eta>0$ instead.
Its lattice separation at level $m$ is at least $\eta/a_m$.
The chosen schedule satisfies $a_mB_m^2\to0$.
Here is the quantitative reason for this limit.
Its support-scale estimate eventually gives $B_{m+1}^2\le2B_m^2$,
whereas the block lengths satisfy $L_m\ge3$ and
$a_{m+1}=a_m/L_m$. Thus $a_{m+1}B_{m+1}^2\le
(2/3)a_mB_m^2$ on a tail; iteration proves the limit.
Choose a common root after all entry levels and after
$a_mB_m^2\le\eta$ holds on that tail. The origins are then
separated by $B_m^2$ at every later level. They are distinct at the
root, so occupancy there is one, within the purchased occupancy
schedule. This verifies all geometric conditions for the preceding
Cauchy argument.

Finally apply the all-period finite-difference bounds of
Section~\ref{sec:ex-fullmesh-fixedbounds}. They give local boundedness and an eventual common modulus on
$\mathrm{NoHit}_p$, the open set of configurations distinct modulo the
physical period. The rounded lattice readouts need not themselves be
continuous; their mesh errors tend to zero. The explicit local bound is
$C(|x-y|+2d a_N)$ for a fixed finite coordinate dimension $d$,
so one first chooses $|x-y|$ small and then the mesh tail. Cauchy convergence on the dense ancestral set in this domain extends
to locally uniform convergence there: cover a compact set by finitely
many neighborhoods with a common equicontinuity modulus, choose an
ancestral point in each, and use the triangle inequality between each
point and its chosen center. The corresponding higher difference
bounds give the smooth limiting correlations by the argument of
Section~\ref{ex-smooth-limits}. Restriction to the ordered Euclidean
configuration domain gives the stated position-space ultraviolet
limit.

In applying this result to the selected construction, the cutoff
parameter is $N$ itself and every branch is the history chosen by the
same shooting construction. The alignment identities identify its
paths with the common horizon; the selected transfer maps and their
based local sections identify all of the $\Phi$ and $\omega$ above.
The common centered service supplies the root allowances and their
positive tails, the transfer ledgers, and the scale estimates.
These identifications are part of the common-family assembly: a
service on a different history would not establish the result here.
